% !TeX root = ../main.tex

\chapter{Chaos Engineering e Test di Fallimento}
\label{cap:chaos_engineering}

Per validare empiricamente i meccanismi di reliability descritti nel Capitolo~\ref{cap:reliability_scalability}, è stata realizzata una suite di chaos engineering (\texttt{ops/chaos\_test\_suite.py}) che inietta in modo controllato guasti rappresentativi sui componenti più critici dell'architettura (database, broker, capacità di flotta). L'obiettivo non è dimostrare l'assenza di guasti, ma misurare la \emph{reazione} del sistema in termini di RTO, perdita di dati e capacità di auto-rimediazione agentica.

\section{Metodologia}
\label{sec:metodologia_chaos}

Ogni esperimento è strutturato secondo il ciclo classico del chaos engineering:

\begin{enumerate}
    \item \textbf{Ipotesi:} definizione del comportamento atteso (SLO da rispettare).
    \item \textbf{Blast radius:} il guasto è confinato a un singolo Deployment nel namespace \texttt{default}; il control-plane non viene toccato.
    \item \textbf{Iniezione:} eseguita via API Kubernetes (\texttt{patch\_namespaced\_deployment\_scale}, \texttt{delete\_namespaced\_pod}) o tramite generatore di carico (\texttt{client\_simulator.py --stress}).
    \item \textbf{Osservazione:} raccolta di metriche su \texttt{chaos\_test\_results.log} (timestamp, RTO, repliche prima/dopo).
    \item \textbf{Abort condition:} timeout massimo di 60\,s per il ripristino, oltre il quale il test è considerato fallito.
\end{enumerate}

\section{Esperimento 1: Primary Database Outage}
\label{sec:esperimento_1_influxdb}

\textbf{Ipotesi.} Lo spegnimento di InfluxDB non deve interrompere il monitoraggio real-time (\textit{Graceful Degradation}, §\ref{sec:graceful_degradation}); la dashboard deve restare reattiva e il sistema deve recuperare entro l'SLO di disponibilità del data layer (§\ref{sec:slo_downtime}).

\textbf{Procedura.} Scaling del Deployment \texttt{influxdb} a 0 repliche, attesa di 10\,s, ripristino a 1 replica e cronometraggio del tempo necessario affinché il Pod torni in stato \texttt{Ready}.

\textbf{Risultato.} \textbf{RTO misurato: 12{,}05\,s} (run del 23/05/2026). Durante l'outage il \texttt{central\_server} ha continuato a servire la dashboard sfruttando lo stato in-memory; le scritture storiche sono state perse, in linea con il pattern di dual-write neutralizzato da \texttt{try/except}.

\section{Esperimento 2: Network Resilience (MQTT Broker Crash)}
\label{sec:esperimento_2_mqtt}

\textbf{Ipotesi.} L'eliminazione forzata del Pod \texttt{mosquitto} non deve compromettere la flotta: Kubernetes deve ricreare il Pod e i client \texttt{paho-mqtt} devono riconnettersi autonomamente tramite \textit{exponential backoff} (§\ref{sec:resilienza_rete}).

\textbf{Procedura.} Deletion del Pod via \texttt{delete\_namespaced\_pod} e attesa del nuovo Pod in stato \texttt{Running}.

\textbf{Risultato.} \textbf{Downtime di scheduling Kubernetes: 0{,}04\,s} (rischedulazione immediata sul nodo). La riconnessione dei droni è gestita dal client MQTT senza logica custom; il broker non perde messaggi grazie alla natura \textit{at-most-once} della telemetria (perdite tollerate, frequenza 2\,s).

\section{Esperimento 3: Load Spike e Scaling Agentico}
\label{sec:esperimento_3_load}

\textbf{Ipotesi.} L'iniezione simultanea di 50 ordini deve essere rilevata dal \texttt{triage\_manager} e tradursi in una richiesta di scaling dell'\texttt{HealthAgent}; sotto la soglia di 6 droni lo scaling avviene in autonomia, oltre tale soglia parte l'escalation HITL (§\ref{sec:hitl_escalation}).

\textbf{Procedura.} Lettura delle repliche correnti del Deployment \texttt{drone-simulator}, esecuzione di \texttt{client\_simulator.py --stress} (durata effettiva $\approx$56\,s), attesa ulteriore di 30\,s per consentire al \texttt{LogisticAIBrain} di completare un ciclo di triage, rilettura delle repliche.

\textbf{Risultato osservato.} Repliche \textbf{prima = 3}, \textbf{dopo = 6}. L'orchestratore agentico ha rilevato il backlog tramite query InfluxDB, invocato il tool MCP \texttt{scale\_drone\_fleet} e raddoppiato la flotta entro $\approx$86\,s end-to-end dalla prima iniezione. Il test conferma l'efficacia dello scaling agentico per shock di carico controllati e mostra che le ottimizzazioni di latenza del ciclo decisionale (\texttt{time.sleep(30)} + inferenza LLM) sono sufficienti per orizzonti dell'ordine del minuto. Per shock sub-minuto resta aperta la traccia (§\ref{sec:lezioni_chaos}) di un HPA Kubernetes nativo come \emph{fast-path} complementare.

\section{Riepilogo Metriche}
\label{sec:riepilogo_chaos}

\begin{table}[ht]
\centering
\begin{tabular}{lll}
\toprule
\textbf{Esperimento} & \textbf{Metrica} & \textbf{Valore} \\
\midrule
InfluxDB outage      & RTO              & 12{,}05\,s \\
MQTT broker crash    & Downtime K8s     & 0{,}04\,s \\
Load spike (50 ord.) & $\Delta$ repliche & 3 $\rightarrow$ 6 ($\approx$86\,s) \\
\bottomrule
\end{tabular}
\caption{Risultati misurati eseguendo \texttt{chaos\_test\_suite.py} (run del 23/05/2026).}
\label{tab:chaos_results}
\end{table}

\section{Lezioni Apprese}
\label{sec:lezioni_chaos}

Gli esperimenti confermano la robustezza dei livelli infrastrutturali (Kubernetes + \texttt{paho-mqtt} + dual-write) e validano lo strato agentico come meccanismo di scaling efficace su orizzonti di $\approx$1\,min. Restano due limiti aperti:
\begin{enumerate}
    \item la latenza del ciclo decisionale ($\approx$30\,s dominati da \texttt{time.sleep(30)} e inferenza LLM) è ancora elevata per spike sub-minuto in cui un HPA Kubernetes nativo, basato su metriche di code/CPU, offrirebbe un \emph{fast-path} complementare;
    \item la suite non misura ancora la \textit{loss rate} di telemetria durante l'outage di InfluxDB né la latenza end-to-end di assegnazione degli ordini sotto stress.
\end{enumerate}
La convivenza HPA + agente (HPA come reattore rapido, agente come pianificatore strategico) rientra tra gli sviluppi futuri discussi nel Capitolo~\ref{cap:conclusioni}.
